\documentclass[11pt]{article}
\usepackage{amsmath,amssymb,amsthm,booktabs,hyperref}
\newtheorem{theorem}{Theorem}
\newtheorem{lemma}{Lemma}
\newcommand{\KL}{D_{\mathrm{KL}}}
\title{A tilted gate-elimination bound for Boolean MCSP}
\author{}\date{}
\begin{document}
\maketitle
\begin{abstract}
At the fixed threshold $\lfloor .99435\,2^n/n\rfloor$, deterministic full-basis circuits for the minimum circuit size problem require more than $2.5226\,2^n$ gates. A promise with thresholds $.993\,2^n/n$ and $.99435\,2^n/n$ still requires more than $2.5168\,2^n$ gates. We combine a one-sided affine hitting property of small truth tables with the circuit census. Tilting the branch probabilities in Chen and Kabanets's elimination process improves circuit shrinkage, while a relative entropy bound guarantees a rejecting point on one of the resulting large affine flats.
\end{abstract}

Write $N=2^n$. The inner size $C(x)$ is the minimum number of fan-in-two Boolean gates needed to compute the $N$-bit table $x$ from an $n$-bit address. The outer circuit has the $N$ table entries as inputs and uses the full binary basis with unrestricted fanout. All numerical bounds below hold for sufficiently large $n$.

For $0<c_1\le c_2<1$, the promise requires acceptance if $C(x)\le\lfloor c_1N/n\rfloor$ and rejection if $C(x)>\lfloor c_2N/n\rfloor$. Put $a=1-c_1$ and $b=1-c_2$. Two facts about MCSP suffice. First, the sharp circuit census bounds the uniform density of the entire accepting set by $2^{-(b-o(1))N}$. Second, every affine flat on which the outer circuit rejects has dimension at most $(a+o(1))N$. Indeed, the dense affine interpolation theorem of the companion manuscript \cite{local} places a table of size at most $(r+o(N))/n$ in every affine flat of codimension $r=\Theta(N)$. A larger flat contains a subflat of any smaller fixed positive dimension. Both applications can thus use a fixed positive margin before passing to a limit.

\begin{theorem}
Choose $0<\eta<1/2$, $3/4\le q<1$ and set
\[
 d=1-H_2(1/2+\eta),\qquad
 D=q\log_2(4q/3)+(1-q)\log_2(4(1-q)).
\]
If $d+(D-2d)_+a/q<b$, every deterministic circuit for this promise has size at least
\[
 \left[2.5+\eta-\max\{3.5+\eta,(3+2\eta)/q\}\,a-o(1)\right]N .
\]
\end{theorem}

\begin{proof}
Use the simplified-circuit restrictions of Chen and Kabanets \cite[Section 4.1, Lemma 4.1]{ck}. Their potential $M$ counts gates and live input variables. An ordinary restriction branches on one nonconstant affine form. In each branch the potential drops by at least three, and in at least one it drops by four. The remaining case uses a bottom AND-type gate with two private raw inputs, each of fanout one. Either gate output eliminates its two inputs, the gate, and a successor. A circuit of potential at most four can be settled by a bounded number of additional equations. A terminal affine function needs at most one equation. All parent restrictions and substitutions are affine when lifted to the original parameters, except for the AND-type step which we handle by sampling its two inputs explicitly.

In an ordinary step choose a branch giving the guaranteed larger decrement with probability $p=1/2+\eta$. At a private AND-type gate, three input pairs give its common output. Choose that output with probability $q$ and sample its three pairs uniformly, or take the exceptional pair otherwise. Tilt only the first $J=\lceil(a+\xi)N/q\rceil$ private gates, using fair raw bits thereafter. Here $\xi>0$ is sufficiently small. Complete any unused parameters uniformly.

If all ordinary branches and private inputs were fair, the resulting point would be uniform. In affine coordinates each chosen ordinary form is a fresh free bit, and each private pair comprises two free bits. One can equivalently sample newly unused parameters immediately and retain their values for the final completion. Hence the sum of the number of ordinary branches and twice the number of private pairs is at most $N$. The chain rule for relative entropy, followed by data processing, bounds the divergence of the actual point law $P$ from uniform $U$ by
\[
 \KL(P\Vert U)\le Nd+(D-2d)_+J.
\]
All relative entropies here are in bits. For the accepting set $E$, binary data processing gives $\KL(P\Vert U)\ge P(E)\log_2(1/U(E))-1$. The assumed strict inequality allows $\xi>0$ such that the sampled point is rejected with probability bounded below by a positive constant.

Let $T$ be the number of nonterminal steps, including $W$ private gates. Truncate each potential decrement at four, pad to $N$ steps by zeros, and center at its conditional mean. Azuma's inequality yields with probability $1-o(1)$
\[
 M_{\rm initial}\ge A T+(4-A)W-o(N),\qquad A=3.5+\eta .
\]
Indeed the conditional mean is at least $A$ at ordinary steps and four at private ones. With the same probability, the number of common outcomes among the first $\min(W,J)$ private gates is $q\min(W,J)+o(N)$, by the corresponding bounded martingale.

Construct an affine flat through the sampled point on which the output stays fixed. Keep one affine equation per ordinary branch. At a private gate with the common output, fix just one input literal that forces this output at the sampled pair. Fix both inputs in the exceptional case. A restored direction is independent of all later choices: in the current affine coordinates the two raw inputs occur nowhere else, and after this gate's output is fixed the reduced circuit is independent of both. Future substitutions involve only the remaining live parameters. Pulling the free direction and the new coordinates back along the injective affine parametrization preserves their independence. Unused coordinates remain free. Thus if $W\ge J$ the flat has dimension at least $(a+\xi)N-o(N)$. Otherwise its dimension $L$ satisfies
\[
 L\ge qW-o(N),\qquad L\ge N-T-(1-q)W-o(N).
\]
The bounded number of terminal equations is absorbed by the remainder.

With positive probability the point is rejected and both concentration estimates hold. The whole flat then rejects and satisfies $L\le(a+o(1))N$, excluding $W\ge J$. Hence $W\le (a/q+o(1))N$. If $G$ is the gate count, substitution gives
\[
 G+N\ge A(N-L)+(4-2A+Aq)W-o(N).
\]
According to the sign of the coefficient of $W$, use either zero or the upper bound on $W$ in this lower bound. The result is the claimed estimate.
\end{proof}

\paragraph{Certified constants.}
Take $q=6/7$, giving $D=q-\log_2(7/4)<1/20$. For transparency,
$\ln 2>6931/10000$, $\ln 2<6932/10000$ and $\ln(7/4)>5596/10000$ follow by inserting $y=1/3$ and $y=3/11$ into
\[
 2\sum_{j=0}^{4}\frac{y^{2j+1}}{2j+1}\le
 \ln\frac{1+y}{1-y}
 \le 2\sum_{j=0}^{4}\frac{y^{2j+1}}{2j+1}
       +\frac{2y^{11}}{11(1-y^2)} .
\]
The indicated bounds imply $\log_2(7/4)>6/7-1/20$. For $x=2\eta$ the entropy expansion gives
\[
 d=\frac1{\ln 2}\sum_{j\ge1}\frac{x^{2j}}{(2j-1)2j}
 \le u:=\frac{x^2}{2(1-x^2)(6931/10000)} .
\]
Both rows below have $2a/q<1$. Their entropy cost is at most
$\max\{u,(1-2a/q)u+a/(20q)\}$.
The last two columns are strict outward rational bounds obtained by substituting the displayed terminating decimals as fractions.
\[
\begin{array}{ccccc}\toprule
 \eta & a & b & \text{entropy cost below} & \text{gate coefficient above}\\\midrule
 .043 & .00565 & .00565 & .005634 & 2.52265\\
 .042 & .007 & .00565 & .005451 & 2.51681\\\bottomrule
\end{array}
\]
This yields the strict constants in the abstract after absorbing $o(1)$. A coin-fixed circuit of a two-sided bounded-error algorithm need not have monochromatic rejecting flats disjoint from the easy set. The proof uses exact correctness on both parts of the promise.
\section*{Exponentially accurate randomized circuits}
The same gate coefficient holds at two-sided pointwise error
$2^{-eN}$ whenever
\[
 e>d+(D-2d)_+a/q.
\]
As in the theorem, assume $\kappa=d+(D-2d)_+a/q<b$. This gives,
for example, the strict $2.5226N$ bound at the first threshold in
the table for every fixed $e\ge .00564$. The same exponent controls uniform false positives and false negatives
on the common easy distribution.

To prove this claim, fix a gate coefficient strictly below the
theorem's bound. Choose constants $a<a'<a''$ sufficiently close to
$a$. Use $a''$ in the twig cutoff and write
$\kappa''=d+(D-2d)_+a''/q$. All strict margins, including
$b>\kappa''$ and $e>\kappa''$, persist. For any
deterministic circuit respecting the new gate bound, the two
concentration estimates in the proof imply that the maximum restored
dimension is at least $a'N$ with tilted probability $1-o(1)$.
This inference only uses the gate count. In particular it holds for
circuits with arbitrary erroneous labels. When the cutoff is not
reached the same algebra applies with $a'$ in place of $a$, since
otherwise the size lower bound with $a'$ would contradict our strict
gate margin.

Consider the generalized terminal path that records each affine
answer and each twig \emph{output}. Choose simplifications from these
answers alone. The likelihood ratio relative to uniform is the
product of the ordinary branch likelihood ratios and the factors
$4q/3$ or $4(1-q)$ at tilted twigs. Its logarithm has bounded
increments, with conditional means $d$ and $D$, respectively.
The pathwise bound on their sum is at most $\kappa''N+O(1)$.
Center the increments, pad to $N$ steps, and apply Azuma's inequality.
Except on a set of tilted mass $o(1)$, the logarithm is at most
$\kappa''N+o(N)$. Thus the set of generalized paths of restored
dimension at least $a'N$ has uniform mass at least
$2^{-(\kappa''+o(1))N}$.

There is an exact fractional representation of each such generalized
path by affine certificates. In a common twig let $H$ be its three
possible pairs. Any two distinct binary points constitute an affine
line. Each point of $H$ belongs to two of its three two-point lines, so
\[
 \mathbf 1_H=\tfrac12\sum_{\ell\subset H,\ |\ell|=2}\mathbf 1_\ell .
\]
The third line, in addition to those that fix one forcing literal,
fixes the parity of the two bits. For each common twig choose one
of these three line equations and retain both literal equations at
exceptional twigs. Tensoring the displayed identity over the $m$
common twigs expresses the path indicator as $2^{-m}$ times the sum
of the $3^m$ resulting certificate indicators. The coordinates used
in each twig are free parameters at that stage, both removed before
later substitutions, so their affine pullbacks give one independent
restored direction per common twig for every choice of lines. In
particular every certificate has the full restored dimension of its
generalized path. This identity permits overlap with its exact
pointwise multiplicity.

Use the dense affine law of the companion manuscript \cite{local}
with threshold $c_1N/n$ and rank slightly above $(1-a')N$.
Its mass in every constituent certificate of a large-dimension path
is at least $1-o(1)$ times the uniform mass there. The positive
weighted identity transfers the same inequality to the entire
generalized path. Let $u$ denote uniform false-positive mass
on inputs of complexity above $c_2N/n$, and let $v$ denote the
false-negative mass under this common easy law. Large-dimension
paths labeled one have uniform mass at most $u+2^{-(b-o(1))N}$
by the census. Those labeled zero have uniform mass at most
$(1+o(1))v$ by fractional domination and disjointness of the
generalized paths. Hence
\[
 u+2^{-(b-o(1))N}+(1+o(1))v
 \ \ge\ 2^{-(\kappa''+o(1))N}.
\]
All asymptotic error bounds are uniform in the fixed circuit.
Averaging over independent circuit coins now gives the claimed
two-sided confidence statement.

\paragraph{A linear dependence on the density deficit at bounded error.}
For fixed $0<c<1$ and threshold $\lfloor cN/n\rfloor$, a bounded-error
outer circuit requires at least $(4c-3/2-o(1))N$ gates. More generally
this conclusion uses the completeness density $c=c_1$ in the promise,
and any fixed soundness density $c_2<1$. Here the two ensembles give
asymptotically identical output distributions for every smaller gate
coefficient.

Indeed use ordinary fair branches and fair raw twig pairs, so the point
law is uniform and every twig is common with probability $3/4$.
The concentration estimates in the proof give, with probability
$1-o(1)$,
\[
 M_{\rm initial}\ge\tfrac72 T+\tfrac12 W-o(N),\qquad
 L\ge \tfrac34 W-o(N),\qquad L\ge N-T-\tfrac14 W-o(N).
\]
For any fixed $a'>1-c$, if $L\le a'N$, substitution yields
$M_{\rm initial}\ge (7/2-4a')N-o(N)$. Since $M_{\rm initial}\le G+N$,
a gate coefficient strictly below $5/2-4(1-c)$ therefore makes
$L>a'N$ typical for some $a'>1-c$. The fractional identity gives
relative domination on every such generalized path. Their uniform
mass is $1-o(1)$, so summing the smaller path probabilities gives
$1-o(1)$ as well. The resulting total variation bound holds for
each fixed circuit and after averaging its coins. Pointwise error
bounded below $1/2$ by a fixed constant would require a positive
distinguishing advantage, since the uniform table is a required no
input with probability $1-o(1)$.


\begin{thebibliography}{9}
\bibitem{ck} R. Chen and V. Kabanets. \emph{Correlation bounds and \#SAT algorithms for small linear-size circuits}. \href{https://www2.cs.sfu.ca/~kabanets/papers/linsize-eccc.pdf}{Author manuscript}, Section 4.
\bibitem{local} Companion manuscript. \emph{Satisfying parity constraints on a Boolean truth table with a small circuit}. Dense local affine interpolation theorem and circuit census. See the accompanying \textnormal{papers/affine/main.tex}.
\end{thebibliography}
\end{document}
